% GENERATED FILE. Edit canonical lessons, then run npm run generate:books.
% canonical-source-hash: fnv1a64:135cfb431dda72e5
% canonical-lessons: PA-C02-naam, PA-C02-mera, PA-C02-hai, PA-C02-mera-naam-hai, PA-C02-tu-tusi, PA-C02-ki, PA-S02-mamma-rara-lava, PA-C02-tuhada-naam-ki-hai, PA-C02-khushi, PA-S02-sassa-tatta-sihari, PA-C02-practice

\chapter{Introducing Yourself}
\label{ch:introductions}
\hlchaptermodality{2}

\begin{chapteropening}
\textbf{By the end of this chapter:} \emph{I can introduce myself in Punjabi, give my name, and understand the other person's name.}

Complete separate oral listening and speaking checks for a respectful Punjabi introduction. Independent Gurmukhi writing remains explicitly blocked: the chapter's two recognition sessions keep the model visible and are scored as recognition only, and romanization earns no script credit.
\end{chapteropening}

\section[nāṁ]{\pa{ਨਾਂ} (nāṁ) — \textquotedblleft{}name,\textquotedblright{} a word older than Europe}

\label{lesson:PA-C02-naam}



\begin{quote}
The first word of the introduction — and one of the clearest bridges between India and Europe.
\end{quote}

\subsection*{The letters in this word}
\textbf{\pa{ਨ}} (\emph{na}) + \textbf{\pa{ਾ}} (the \emph{kannā}, long-\emph{ā} sign) + \textbf{\pa{ਂ}} (the \emph{bindi}, a nasal dot) $\to$ \textbf{\pa{ਨਾਂ}} (\emph{nāṁ}). A single nasalised syllable.

\begin{cousinweb}
\textbf{\pa{ਨਾਂ}} (\emph{nāṁ}, \textquotedblleft{}name\textquotedblright{}) $\leftarrow$ Sanskrit \emph{nāman} (a Sanskrit word, so it is written here the way Sanskrit is written in this book — in roman letters, not in Gurmukhi), from Proto-Indo-European \textbf{*h\textsubscript{3}nómn} (with a syllabic final \emph{n}) — the \emph{same} ancient root as English \textbf{name}, Latin \emph{nōmen} (English \textbf{noun}, \textbf{nom}inate), and Greek \emph{onoma}. So Punjabi \emph{nāṁ} and English \emph{name} are cousins from before either language existed. (The Hindi form is \emph{nām}; Punjabi nasalises the ending to \emph{nāṁ}.)
\end{cousinweb}

\subsection*{Your turn}
Say these aloud:

\begin{itemize}
\raggedright
  \item \textquotedblleft{}nāṁ\textquotedblright{}
  \item the family — \emph{nāṁ} / English \emph{name} / Latin \emph{nōmen} / Greek \emph{onoma}
\end{itemize}

\begin{tcolorbox}[breakable,colback=teal!4,colframe=teal!35!black,title={Before you move on}]
What Sanskrit and PIE roots is \emph{nāṁ} from, and its English cousins? (\emph{Nāman} / \emph{*h\textsubscript{3}nómn}, with a syllabic final \emph{n}; \textbf{name}, \textbf{noun}, \textbf{nominate}.)
\end{tcolorbox}
\section[merā]{\pa{ਮੇਰਾ} (merā) — \textquotedblleft{}my\textquotedblright{}}

\label{lesson:PA-C02-mera}



\begin{quote}
The little word that makes the name \emph{yours}.
\end{quote}

\subsection*{The letters in this word}
\textbf{\pa{ਮ}} (\emph{ma}) + \textbf{\pa{ੇ}} (the \emph{lāvā̃}, \emph{e}-sign) $\to$ \textbf{\pa{ਮੇ}} (\emph{me}); \textbf{\pa{ਰ}} (\emph{ra}) + \textbf{\pa{ਾ}} (long \emph{ā}) $\to$ \textbf{\pa{ਰਾ}} (\emph{rā}). Read \textbf{\pa{ਮੇ}·\pa{ਰਾ}} $\to$ \emph{merā}.

\begin{cousinweb}
\textbf{\pa{ਮੇਰਾ}} (\emph{merā}, \textquotedblleft{}my\textquotedblright{}) is on the first-person root \textbf{ma-} — the same \emph{ma-/me-} behind English \textbf{me}, \textbf{my}, \textbf{mine} and Latin \emph{meus}, and the same root as Hindi \emph{merā}. The Punjabi \textquotedblleft{}my\textquotedblright{} and the English \textquotedblleft{}my\textquotedblright{} are, at the root, one word.
\end{cousinweb}

\begin{grammarlens}[title={Grammar Lens: \textquotedblleft{}my\textquotedblright{} agrees with the noun}]
Punjabi possessives change to match what is owned: \textbf{\pa{ਮੇਰਾ}} (\emph{merā}, masculine), \emph{merī} (feminine), \emph{mere} (plural). \emph{nāṁ} (\textquotedblleft{}name\textquotedblright{}) is masculine, so it takes \emph{merā}: \emph{merā nāṁ}. (You will meet \emph{merī} the first time you own a feminine noun.)
\end{grammarlens}

\subsection*{Your turn}
Say these aloud:

\begin{itemize}
\raggedright
  \item \textquotedblleft{}merā\textquotedblright{}
  \item the m- family — \emph{merā} (my), English \emph{me}, \emph{my}
  \item the feminine form (\emph{merī})
\end{itemize}

\begin{tcolorbox}[breakable,colback=teal!4,colframe=teal!35!black,title={Before you move on}]
What root links \emph{merā} and English \emph{my}, and what is its feminine form? (The first-person \emph{ma-}; \emph{merī}.)
\end{tcolorbox}
\section[hai]{\pa{ਹੈ} (hai) — \textquotedblleft{}is,\textquotedblright{} cousin of English \emph{is}}

\label{lesson:PA-C02-hai}



\begin{quote}
The little verb \textquotedblleft{}is\textquotedblright{} — the same \textquotedblleft{}is\textquotedblright{} spoken from Ireland to India.
\end{quote}

\subsection*{The letters in this word}
\textbf{\pa{ਹ}} (\emph{ha}) + \textbf{\pa{ੈ}} (the \emph{dulāvā̃}, \emph{ai}-sign) $\to$ \textbf{\pa{ਹੈ}} (\emph{hai}). One syllable.

\begin{cousinweb}
\textbf{\pa{ਹੈ}} (\emph{hai}, \textquotedblleft{}is\textquotedblright{}) $\leftarrow$ Sanskrit \emph{asti} (\textquotedblleft{}is\textquotedblright{}), from Proto-Indo-European \textbf{*h\textsubscript{1}es-} — the root of \textquotedblleft{}to be\textquotedblright{} across the whole family: English \textbf{is}, German \emph{ist}, Latin \emph{est}, Spanish \emph{es}. Punjabi \emph{hai} is the Indian cousin of the same tiny verb.
\end{cousinweb}

\begin{grammarlens}[title={Grammar Lens: the verb comes last}]
Punjabi is \textbf{subject–object–verb}: \textquotedblleft{}my name Arun \textbf{is}.\textquotedblright{} \emph{hai} sits at the \emph{end} of the sentence, where English puts \textquotedblleft{}is\textquotedblright{} in the middle. Bank the shape — in Punjabi the verb waits until the end.
\end{grammarlens}

\subsection*{Your turn}
Say these aloud:

\begin{itemize}
\raggedright
  \item \textquotedblleft{}hai\textquotedblright{}
  \item the family — \emph{hai} / \emph{is} / \emph{ist} / \emph{est} / \emph{es}
  \item \textquotedblleft{}merā nāṁ Arun hai\textquotedblright{} — note \emph{hai} at the end
\end{itemize}

\begin{tcolorbox}[breakable,colback=teal!4,colframe=teal!35!black,title={Before you move on}]
What root is \emph{hai} from, and where does the verb sit in a Punjabi sentence? (Sanskrit \emph{asti} / PIE \emph{*h\textsubscript{1}es-}; at the end.)
\end{tcolorbox}
\section[merā nāṁ … hai]{\pa{ਮੇਰਾ} \pa{ਨਾਂ} … \pa{ਹੈ} (merā nāṁ … hai) — \textquotedblleft{}my name is…\textquotedblright{}}

\label{lesson:PA-C02-mera-naam-hai}



\begin{quote}
Three atoms you now own, assembled into a first sentence.
\end{quote}

\subsection*{The exchange — assembled}
\textbf{\pa{ਮੇਰਾ}} (my) + \textbf{\pa{ਨਾਂ}} (name) + your name + \textbf{\pa{ਹੈ}} (is) = \textbf{merā nāṁ … hai}, \textquotedblleft{}my name is…\textquotedblright{} \emph{Merā nāṁ Arun hai.} $\to$ \textquotedblleft{}My name is Arun.\textquotedblright{}

\begin{grammarlens}[title={Grammar Lens: the verb goes last}]
Punjabi is a \textbf{subject–object–verb} language — literally \textquotedblleft{}my name Arun \textbf{is}.\textquotedblright{} The \emph{hai} sits at the end, where English puts \textquotedblleft{}is\textquotedblright{} in the middle. Bank this shape: in Punjabi, the verb waits until the end of the sentence.
\end{grammarlens}

\subsection*{Your turn}
Say these aloud:

\begin{itemize}
\raggedright
  \item \textquotedblleft{}merā nāṁ … hai\textquotedblright{} with your name
  \item notice the verb \emph{hai} comes last
\end{itemize}

\begin{tcolorbox}[breakable,colback=teal!4,colframe=teal!35!black,title={Before you move on}]
Give your name in Punjabi, and say where the verb sits. (\emph{Merā nāṁ … hai}; the verb is last.)
\end{tcolorbox}
\section[tū̃ / tusī̃]{\pa{ਤੂੰ} / \pa{ਤੁਸੀਂ} (tū̃ / tusī̃) — \textquotedblleft{}you,\textquotedblright{} familiar and respectful}

\label{lesson:PA-C02-tu-tusi}



\begin{quote}
To ask a name you need \textquotedblleft{}you\textquotedblright{} — and Punjabi grades respect the way French does.
\end{quote}

\subsection*{The letters in this word}
\textbf{\pa{ਤੂੰ}}: \textbf{\pa{ਤ}} (\emph{ta}) + \textbf{\pa{ੂ}} (long \emph{ū}) + \textbf{\pa{ਂ}} (nasal \emph{bindi}) $\to$ \emph{tū̃}. \textbf{\pa{ਤੁਸੀਂ}}: \emph{tu} + \emph{sī̃}, a longer, plural-shaped word.

\begin{cousinweb}
\textbf{\pa{ਤੂੰ}} (\emph{tū̃}, familiar \textquotedblleft{}you\textquotedblright{}) $\leftarrow$ Sanskrit \emph{tvam}, from PIE \textbf{*tū} — the same root as English archaic \textbf{thou}, Latin \emph{tū}, French \emph{tu}. \textbf{\pa{ਤੁਸੀਂ}} (\emph{tusī̃}, respectful \textquotedblleft{}you\textquotedblright{}) is grammatically plural — the same courtesy-by-plural used by French \emph{vous} and Hindi \emph{āp}.
\end{cousinweb}

\begin{grammarlens}[title={Grammar Lens: two levels of \textquotedblleft{}you\textquotedblright{}}]
\begin{itemize}
\raggedright
  \item \textbf{\pa{ਤੂੰ}} (\emph{tū̃}) — familiar: friends, children, intimates.
  \item \textbf{\pa{ਤੁਸੀਂ}} (\emph{tusī̃}) — respectful: elders, strangers, politeness (and true plural, \textquotedblleft{}you all\textquotedblright{}).
\end{itemize}

For a first meeting, always \textbf{\pa{ਤੁਸੀਂ}}. Respect takes the plural verb, exactly as in the Romance and other Indo-Aryan tracks.
\end{grammarlens}

\subsection*{Your turn}
Say these aloud:

\begin{itemize}
\raggedright
  \item \textquotedblleft{}tū̃\textquotedblright{} (familiar) and \textquotedblleft{}tusī̃\textquotedblright{} (respectful)
  \item the English cousin of \emph{tū̃} (\textbf{thou})
  \item which you use on first meeting (\emph{tusī̃})
\end{itemize}

\begin{tcolorbox}[breakable,colback=teal!4,colframe=teal!35!black,title={Before you move on}]
Which \textquotedblleft{}you\textquotedblright{} is respectful, and what English word is \emph{tū̃} cousin to? (\emph{Tusī̃}; archaic \textbf{thou}, from PIE \emph{*tū}.)
\end{tcolorbox}
\section[kī]{\pa{ਕੀ} (kī) — \textquotedblleft{}what\textquotedblright{}}

\label{lesson:PA-C02-ki}



\begin{quote}
The question-word — and it starts with the \textbf{k-} that heads nearly every Punjabi question.
\end{quote}

\subsection*{The letters in this word}
\textbf{\pa{ਕ}} (\emph{ka}) + \textbf{\pa{ੀ}} (the \emph{bihārī}, long-\emph{ī} sign) $\to$ \textbf{\pa{ਕੀ}} (\emph{kī}).

\begin{cousinweb}
\textbf{\pa{ਕੀ}} (\emph{kī}, \textquotedblleft{}what\textquotedblright{}) $\leftarrow$ Sanskrit \emph{kim} / the interrogative stem \textbf{ka-}, from PIE \textbf{*k\textsuperscript{w}o-} — the same question-root behind English \textbf{wh-} words (\textbf{what}, \textbf{who}, \textbf{which}) and Latin \emph{qu-}. Punjabi's question-words nearly all begin with this \textbf{k-}: \emph{kī} (what), \emph{kauṇ} (who), \emph{kad} (when), \emph{kithē} (where), \emph{kivēṁ} (how). Learn the \emph{k-} and you can hear a question coming.
\end{cousinweb}

\subsection*{Your turn}
Say these aloud:

\begin{itemize}
\raggedright
  \item \textquotedblleft{}kī\textquotedblright{}
  \item the k- question family — \emph{kī, kauṇ, kad, kithē, kivēṁ}
  \item the English cousins (\emph{what, who, which})
\end{itemize}

\begin{tcolorbox}[breakable,colback=teal!4,colframe=teal!35!black,title={Before you move on}]
What root do Punjabi's question-words share, and its English family? (The interrogative \emph{k-}, PIE \emph{*k\textsuperscript{w}o-}; the \emph{wh-} words.)
\end{tcolorbox}
\section[ma · ra · -e]{\pa{ਮ} · \pa{ਰ} · \pa{ੇ} — three pieces met inside a word you already say}

\label{lesson:PA-S02-mamma-rara-lava}



\begin{quote}
You are about to say \emph{merā nāṁ} out loud for a whole chapter. Three of the marks that spell it stop being decoration now.

First, say the last things you learned once more, out loud and from memory:

\begin{itemize}
\raggedright
  \item \emph{Say it:} \emph{dhannavād} · \emph{shukrīā}
\end{itemize}
\end{quote}

\subsection*{Script you'll notice: \pa{ਮ} · \pa{ਰ} · \pa{ੇ}}
\begin{quote}
\pa{ਮ}
\end{quote}

\textbf{mammā} — \emph{ma}. A consonant. In Gurmukhi a bare consonant already carries an \emph{a}, so this is \emph{ma}, never a bare \emph{m}.

\begin{quote}
\pa{ਰ}
\end{quote}

\textbf{rārā} — \emph{ra}. A consonant, tapped once with the tip of the tongue — closer to Spanish \emph{pero} than to English \emph{r}.

\begin{quote}
\pa{ੇ}
\end{quote}

\textbf{lāvā̃} — \emph{-e}. Not a letter. A \textbf{mātrā} — a vowel sign — and it rides \textbf{on the head-line above} the consonant it changes.

Now put them where you have already heard them:

\begin{quote}
\textbf{\pa{ਮ}} + \textbf{\pa{ੇ}} $\to$ \textbf{\pa{ਮੇ}} (\emph{me});  \textbf{\pa{ਰ}} + \textbf{\pa{ਾ}} $\to$ \textbf{\pa{ਰਾ}} (\emph{rā})
\end{quote}

>

\begin{quote}
\textbf{\pa{ਮੇਰਾ}} — \emph{merā} — \textbf{my}
\end{quote}

You are not being asked to write any of this. You are being asked to \textbf{recognise} it the next time it goes past.

\subsection*{Writing: copy what you can see}
Keep \pa{ਮ} · \pa{ਰ} · \pa{ੇ} in front of you and do not cover any of them. Follow each shape once with a finger, then copy it once with a pencil, larger than it is printed. Say its sound as your hand stops.

\begin{quote}
This is \textbf{observe-and-trace with the model visible}. Nothing in these chapters asks you to write a Punjabi word from memory; that ladder starts later, and it starts from these same shapes.
\end{quote}

\subsection*{Your turn}
\begin{itemize}
\raggedright
  \item \emph{Look:} at \textbf{\pa{ਮੇਰਾ}} and point to each of today's pieces in turn
  \item \emph{Say it:} \emph{ma}, \emph{ra}, \emph{-e}, then the whole word — \emph{merā}
  \item \emph{Look:} back at the headword at the top of any earlier lesson in this chapter and find one shape you now know
\end{itemize}

\begin{tcolorbox}[breakable,colback=teal!4,colframe=teal!35!black,title={Before you move on}]
Name today's pieces: \emph{mammā}, \emph{rārā}, \emph{lāvā̃}. Which of them is a \textbf{mātrā} or a mark rather than a letter, and what does \textbf{\pa{ਮੇਰਾ}} say? (\emph{merā}, “my”.)

Source: \href{https://www.unicode.org/charts/PDF/U0A00.pdf}{Unicode Gurmukhi chart}.
\end{tcolorbox}
\section[tuhāḍā nāṁ kī hai?]{\pa{ਤੁਹਾਡਾ} \pa{ਨਾਂ} \pa{ਕੀ} \pa{ਹੈ}? (tuhāḍā nāṁ kī hai?) — \textquotedblleft{}what's your name?\textquotedblright{}}

\label{lesson:PA-C02-tuhada-naam-ki-hai}



\begin{quote}
Turn the introduction around — ask, instead of tell.
\end{quote}

\subsection*{The exchange — assembled}
\textbf{\pa{ਤੁਹਾਡਾ}} (\emph{tuhāḍā}, \textquotedblleft{}your,\textquotedblright{} the respectful possessive from \emph{tusī̃}) + \textbf{\pa{ਨਾਂ}} (name) + \textbf{\pa{ਕੀ}} (what) + \textbf{\pa{ਹੈ}} (is) $\to$ literally \textquotedblleft{}\textbf{your name what is?}\textquotedblright{} Punjabi keeps the verb last again, and simply drops the question-word \emph{kī} into the object slot.

\begin{grammarlens}[title={Grammar Lens: answering}]
Answer it with the sentence you built: \emph{Merā nāṁ Arun hai.} The familiar version swaps \emph{tuhāḍā} for \textbf{\pa{ਤੇਰਾ}} (\emph{terā}, \textquotedblleft{}your,\textquotedblright{} from \emph{tū̃}): \emph{terā nāṁ kī hai?}
\end{grammarlens}

\subsection*{Your turn}
Say these aloud:

\begin{itemize}
\raggedright
  \item the respectful question — \emph{tuhāḍā nāṁ kī hai?}
  \item answer it — \emph{merā nāṁ … hai}
  \item the familiar \textquotedblleft{}your\textquotedblright{} (\emph{terā})
\end{itemize}

\begin{tcolorbox}[breakable,colback=teal!4,colframe=teal!35!black,title={Before you move on}]
Ask someone's name respectfully, and note where the verb sits. (\emph{Tuhāḍā nāṁ kī hai?}; the verb is last.)
\end{tcolorbox}
\section[khushī hoī]{\pa{ਖ਼ੁਸ਼ੀ} \pa{ਹੋਈ} (khushī hoī) — \textquotedblleft{}pleased to meet you\textquotedblright{}}

\label{lesson:PA-C02-khushi}



\begin{quote}
The warm close of an introduction — and, like \emph{shukrīā}, it comes from Persian, not Sanskrit.
\end{quote}

\subsection*{The letters in this word}
\textbf{\pa{ਖ਼ੁਸ਼ੀ}} (\emph{khushī}): note \textbf{two} \emph{pair bindis} (dots beneath) — \textbf{\pa{ਖ਼}} (\emph{kh}) and \textbf{\pa{ਸ਼}} (\emph{sh}) — both marking Perso-Arabic sounds, exactly as you saw on \textbf{\pa{ਸ਼}} in \emph{shukrīā} (Chapter 1).

\begin{cousinweb}
The Punjabi \textquotedblleft{}pleased to meet you\textquotedblright{} is \emph{tuhānū̃ mil ke khushī hoī} — literally \textquotedblleft{}\textbf{having met you, happiness happened}.\textquotedblright{} The key word \textbf{\pa{ਖ਼ੁਸ਼ੀ}} (\emph{khushī}, \textquotedblleft{}happiness\textquotedblright{}) is \textbf{not} Sanskrit — it is \textbf{Persian} (\emph{khush}, \textquotedblleft{}happy, pleasant\textquotedblright{}), one of the thousands of Persian words woven through Punjabi. So even introducing yourself draws on Punjabi's \textbf{two vocabularies}: Sanskrit (\emph{nāṁ}, \emph{hai}) and Perso-Arabic (\emph{khushī}) — the thread you first met in Chapter 1's \emph{shukrīā}.
\end{cousinweb}

\subsection*{Your turn}
Say these aloud:

\begin{itemize}
\raggedright
  \item \textquotedblleft{}khushī hoī\textquotedblright{}
  \item which heritage \emph{khushī} belongs to (Persian, not Sanskrit)
  \item what the two \emph{pair bindis} (\pa{ਖ਼}, \pa{ਸ਼}) mark (Perso-Arabic sounds)
\end{itemize}

\begin{tcolorbox}[breakable,colback=teal!4,colframe=teal!35!black,title={Before you move on}]
What does the phrase literally mean, and where does \emph{khushī} come from? (\textquotedblleft{}Having met you, happiness happened\textquotedblright{}; Persian \emph{khush}, \textquotedblleft{}happy\textquotedblright{} — Punjabi's second vocabulary.)
\end{tcolorbox}
\section[sa · ta · ka]{\pa{ਸ} · \pa{ਤ} · \pa{ਕ} — three pieces met inside a word you already say}

\label{lesson:PA-S02-sassa-tatta-sihari}



\begin{quote}
Three more consonants — and with the two you met a moment ago, whole words start falling out.

First, say the last things you learned once more, out loud and from memory:

\begin{itemize}
\raggedright
  \item \emph{Say it:} \emph{hai} · \emph{merā nāṁ … hai} · \emph{tū̃ / tusī̃}
\end{itemize}
\end{quote}

\subsection*{Script you'll notice: \pa{ਸ} · \pa{ਤ} · \pa{ਕ}}
\begin{quote}
\pa{ਸ}
\end{quote}

\textbf{sassā} — \emph{sa}. A consonant, a plain \emph{s}. It is also the letter the whole script is often taught from first.

\begin{quote}
\pa{ਤ}
\end{quote}

\textbf{tattā} — \emph{ta}. A consonant, the \textbf{dental} \emph{t}: tongue on the teeth, softer than any English \emph{t}.

\begin{quote}
\pa{ਕ}
\end{quote}

\textbf{kakkā} — \emph{ka}. A consonant, the plain \emph{k} of \emph{skin} — no puff of breath, unlike the \emph{k} of English \emph{kin}.

Now put them where you have already heard them:

\begin{quote}
\textbf{\pa{ਸ}} + \textbf{\pa{ਾ}} + \textbf{\pa{ਰ}} + \textbf{\pa{ਾ}} $\to$ \textbf{\pa{ਸਾਰਾ}} (\emph{sārā}, “whole”)
\end{quote}

>

\begin{quote}
\textbf{\pa{ਤ}} + \textbf{\pa{ਾ}} + \textbf{\pa{ਰ}} + \textbf{\pa{ਾ}} $\to$ \textbf{\pa{ਤਾਰਾ}} (\emph{tārā}, “star”)
\end{quote}

>

\begin{quote}
\textbf{\pa{ਕ}} + \textbf{\pa{ਰ}} $\to$ \textbf{\pa{ਕਰ}} (\emph{kar}, “do!”)
\end{quote}

>

\begin{quote}
\textbf{\pa{ਤਾਰਾ}} — \emph{tārā} — \textbf{star}
\end{quote}

You are not being asked to write any of this. You are being asked to \textbf{recognise} it the next time it goes past.

\subsection*{Writing: copy what you can see}
Keep \pa{ਸ} · \pa{ਤ} · \pa{ਕ} in front of you and do not cover any of them. Follow each shape once with a finger, then copy it once with a pencil, larger than it is printed. Say its sound as your hand stops.

\begin{quote}
This is \textbf{observe-and-trace with the model visible}. Nothing in these chapters asks you to write a Punjabi word from memory; that ladder starts later, and it starts from these same shapes.
\end{quote}

\subsection*{Your turn}
\begin{itemize}
\raggedright
  \item \emph{Look:} at \textbf{\pa{ਤਾਰਾ}} and point to each of today's pieces in turn
  \item \emph{Say it:} \emph{sa}, \emph{ta}, \emph{ka}, then the whole word — \emph{tārā}
  \item \emph{Look:} back at the headword at the top of any earlier lesson in this chapter and find one shape you now know
\end{itemize}

\begin{tcolorbox}[breakable,colback=teal!4,colframe=teal!35!black,title={Before you move on}]
Name today's pieces: \emph{sassā}, \emph{tattā}, \emph{kakkā}. Which of them is a \textbf{mātrā} or a mark rather than a letter, and what does \textbf{\pa{ਤਾਰਾ}} say? (\emph{tārā}, “star”.)

Source: \href{https://www.unicode.org/charts/PDF/U0A00.pdf}{Unicode Gurmukhi chart}.
\end{tcolorbox}
\section[Practice]{Chapter 2 — The introduction exchange}

\label{lesson:PA-C02-practice}



\begin{quote}
No new words. The whole introduction, assembled from atoms you now own.
\end{quote}

\subsection*{The exchange}
Say these aloud:

\begin{itemize}
\raggedright
  \item \textquotedblleft{}My name is Mira.\textquotedblright{} — \emph{merā nāṁ Mira hai.}
  \item \textquotedblleft{}What's your name?\textquotedblright{} — \emph{tuhāḍā nāṁ kī hai?}
  \item \textquotedblleft{}My name is Arun.\textquotedblright{} — \emph{merā nāṁ Arun hai.}
  \item \textquotedblleft{}Pleased to meet you.\textquotedblright{} — \emph{tuhānū̃ mil ke khushī hoī.}
\end{itemize}

\subsection*{Your turn — the atoms}
Say these aloud:

\begin{itemize}
\raggedright
  \item \textquotedblleft{}name,\textquotedblright{} \textquotedblleft{}my,\textquotedblright{} \textquotedblleft{}is\textquotedblright{} (\emph{nāṁ}, \emph{merā}, \emph{hai})
  \item the two \textquotedblleft{}you\textquotedblright{} levels (\emph{tū̃} / \emph{tusī̃})
  \item \textquotedblleft{}what\textquotedblright{} and where the verb sits (\emph{kī}; the verb is last)
\end{itemize}

\begin{cousinweb}
\begin{itemize}
\raggedright
  \item \emph{nāṁ} $\leftarrow$ Sanskrit \emph{nāman} $\to$ English \textbf{name}
  \item \emph{merā} $\leftarrow$ \emph{ma-} $\to$ English \textbf{my}; \emph{hai} $\leftarrow$ \emph{asti} $\to$ English \textbf{is}
  \item \emph{kī} $\leftarrow$ \emph{ka-} $\to$ English \textbf{what}; \emph{tū̃} $\leftarrow$ \emph{*tū} $\to$ \textbf{thou}
  \item \emph{khushī} $\leftarrow$ \textbf{Persian} — Punjabi's second vocabulary
\end{itemize}
\end{cousinweb}

\begin{tcolorbox}[breakable,colback=teal!4,colframe=teal!35!black,title={Before you move on}]
Independent Gurmukhi reading and writing are \textbf{not scored here}. Chapter 2 still contains untaught forms; \#13068 must build their visible-model, guided-copy, delayed-copy, and dictation runway before either skill can earn credit. A romanized answer never counts as Gurmukhi writing.

Introduce yourself and ask a stranger's name in Punjabi. (\emph{Merā nāṁ … hai. Tuhāḍā nāṁ kī hai?})
\end{tcolorbox}
